\section{Game solutions and matching: proofs}
\label{sec:annex-p18-game-argumentation}

\begin{annexitem}{annex:p18.imputation.core}{thm:p18-imputation-core}{Efficient extension of a blocking coalition}
\begin{proof}
Repeated superadditivity gives
\(v(N)\geq\sum_i v(\{i\})\). Give every player its
singleton value and allocate the nonnegative surplus
to one player. This is an imputation, proving
nonemptiness. Strict improvement on a nonempty
coalition excludes self-domination.

If \(x\succ_K y\), then \(y(K)<x(K)\leq v(K)\),
so a point satisfying every core inequality is
undominated. Conversely suppose \(y(K)<v(K)\)
for an imputation \(y\). The violating coalition
is nonempty because the empty values are zero,
and proper because \(y(N)=v(N)\).
Put \(\delta=v(K)-y(K)>0\), and set
\(x_i=y_i+\delta/|K|\) for \(i\in K\).
These players strictly improve, remain individually
rational and receive total \(v(K)\).
For outsiders first set \(x_j=v(\{j\})\).
Superadditivity gives
\[
                    v(N)\geq v(K)+\sum_{j\notin K}v(\{j\}).
\]
Allocate the remaining nonnegative surplus to
one outsider, which exists because \(K\ne N\).
The resulting \(x\) is individually rational
and efficient and dominates \(y\) through \(K\).
This proves the reverse implication with an
actual feasible imputation.
\end{proof}
\end{annexitem}

\begin{annexitem}{annex:p18.nm.argumentation}{thm:p18-nm-argumentation}{Both cooperative-game correspondences}
\begin{proof}
The internal NM condition is exactly conflict-freeness
in the domination graph, and its external condition
is exactly attack on every outsider. These are
the two stable-extension conditions, proving both
directions on the full imputation set.

An argument is defended by the empty set exactly
when it has no attacker. Here that says its
imputation is undominated, so
\(F_v(\varnothing)={\cal C}(v)\).
This first grounded successor is contained
in the grounded extension. A core point
outside an NM solution would be dominated
by that solution's external condition,
which is impossible. Preferred existence
is the previously proved arbitrary-framework
maximality theorem, applicable to the set
\(I(v)\) regardless of its cardinality.
It supplies a maximal admissible set,
not necessarily a stable one.
\end{proof}
\end{annexitem}

\begin{annexitem}{annex:p18.three.player.solutions}{thm:p18-three-player-solutions}{The two full solution families and their boundaries}
\begin{proof}
Two disjoint coalitions cannot both have at least two
members. If neither does, their values are zero;
if one does, it and their union have value 2
and the other has value zero. This proves
superadditivity, including empty coalitions.

Between distinct vertices of \(S_\triangle\),
one coordinate is equal, one increases and
one decreases. Only a singleton could
strictly improve every member, but singleton
feasibility would require its improved
nonnegative payoff to be at most zero.
Thus internal stability holds. For an
outsider \(y\), at least two coordinates
are less than 1: otherwise two are
at least 1, and nonnegativity with
sum 2 forces exactly a vertex of
\(S_\triangle\). Give those two players
1 each and the third zero. This
vertex strictly improves the two
players and pays their coalition
exactly \(2=v(K)\), proving external
stability with equality permitted
in the feasibility bound.

All points in \(S_c\) have coordinate
3 equal to \(c\), and the first
two sum to \(2-c\). No improving
coalition can contain player 3;
players 1 and 2 cannot both
strictly improve; and singletons
fail feasibility as above.
Thus internal stability holds
throughout \(0\leq c\leq2\).

Suppose \(c<1\) and \(y\notin S_c\).
If \(y_3>c\), put \(\delta=y_3-c\)
and take
\[
                       x=(y_1+\delta/2,y_2+\delta/2,c).
\]
It belongs to \(S_c\) and
dominates \(y\) through \(\{1,2\}\),
whose payoff sum is \(2-c\leq2\).
If \(y_3<c\), at least one
of \(y_1,y_2\) is less than
\(2-c\), since otherwise their
sum is at least \(4-2c>2\).
Choose such an \(i\) and
\(y_i<x_i\leq2-c\), put \(x_3=c\),
and give the other player
\(2-c-x_i\). This belongs to
\(S_c\) and dominates \(y\)
through \(\{i,3\}\), with
coalition sum at most 2.
For \(c=0\) this second
case is empty. Equality
\(y_3=c\) would put \(y\)
inside \(S_c\); thus all
outsiders are covered.

For \(1\leq c\leq2\), the
outsider \(y=(1,1,0)\) cannot
be dominated by \(S_c\).
Both first coordinates of
any \(x\in S_c\) are at
most \(2-c\leq1\), so
only player 3 can improve;
its singleton feasibility
would require \(c\leq0\).
This proves failure also
at the endpoints \(c=1,2\).
Different \(c\)'s give
different nonempty sets,
proving infinitely many
NM solutions.

The three pair inequalities
for a core point would
sum to \(2(y_1+y_2+y_3)=4\geq6\).
Thus the core is empty.
Then \(F_v(\varnothing)=\varnothing\),
so the grounded extension
is empty as well.
Finally take
\[
 x=(1,1,0),\qquad y=(1/2,1/2,1),\qquad z=(0,3/2,1/2).
\]
One has \(x\succ_{\{1,2\}}y\)
and \(y\succ_{\{1,3\}}z\).
Only player 1 improves
from \(z\) to \(x\), and
its singleton feasibility
fails. Thus \(x\) does
not dominate \(z\),
proving nontransitivity.
\end{proof}
\end{annexitem}

\begin{annexitem}{annex:p18.marriage.argumentation}{thm:p18-marriage-argumentation}{Partial functions, totality and every blocking case}
\begin{proof}
Two distinct pairs sharing a participant
are strictly comparable by that person's
total order, so one attacks the other.
Conflict-freeness therefore gives a
partial matching. Conversely disjoint
pairs cannot attack and strict
irreflexivity excludes self-attacks.

Let \(\mu\) be a stable bijection.
For any unchosen pair \((m,w)\),
the two participants do not both
prefer it. Their current partners
are different from the proposed
ones, so strict completeness gives
either \(\mu(m)\succ_m w\) or
\(\mu^{-1}(w)\succ_w m\).
The corresponding chosen incident
pair attacks \((m,w)\).
Thus the graph of \(\mu\)
is a stable extension.

Conversely let \(S\) be stable,
hence a partial matching.
If a man is unmatched,
finite equal cardinalities
give an unmatched woman.
Their unchosen pair has
no chosen incident pair
and therefore no attacker
in \(S\), contradicting
stability. Thus the
matching is a bijection.
A blocking pair for it
could not be attacked
by either chosen incident
pair, and all other
chosen pairs have different
coordinates. This is again
an unattacked outsider.
Hence the bijection is
stable. For \(n=0\)
the sole relation is
the empty bijection,
and the empty framework
has its unique stable
extension empty.
\end{proof}
\end{annexitem}

\begin{annexitem}{annex:p18.deferred.acceptance}{thm:p18-deferred-acceptance}{Deferred acceptance and both optimality directions}
\begin{proof}
Every receiver holds at most one proposer,
and only unheld proposers act, so
held pairs always form a partial
matching. Once a receiver receives
a proposal she never becomes empty;
her held proposer only improves
in her strict order.

There are \(n^2\) possible proposals
and none repeats. If a free
proposer had already proposed to
every receiver, all receivers
would be holding distinct men,
exhausting \(M\), a contradiction.
Thus a permitted proposal exists
whenever a free proposer remains.
The procedure stops within
\(n^2\) proposals with a
perfect matching.

If a man prefers an unchosen
woman to his final partner,
he proposed to her earlier
and was rejected. She
then held someone she
preferred and never worsened
her choice, so she prefers
her final partner to him.
There is no blocking pair.

Induct over rejections to
prove that a rejected pair
belongs to no stable matching.
Suppose \(w\) rejects \(m\)
in favor of \(m'\), and
the claim holds for earlier
rejections. Every woman
\(m'\) prefers to \(w\)
has already rejected him.
In a hypothetical stable
matching containing \((m,w)\),
\(m'\) can have neither
one of those better
partners, by induction,
nor \(w\), who is
assigned to \(m\).
He therefore prefers
\(w\) to his hypothetical
partner, while \(w\)
prefers \(m'\) to \(m\).
They block, a contradiction.
This proves the rejection
claim, starting from the
empty rejection history.

Every partner preferred
by a proposer to his
output has rejected him,
so none is available
to him in any stable
matching. This proves
proposer optimality.
Let \(\mu\) be the output
and \(\nu\) any stable
matching. If a woman
\(w\) strictly preferred
her \(\mu\)-partner \(m\)
to her \(\nu\)-partner,
proposer optimality would
make \(m\) prefer his
distinct \(\mu\)-partner
\(w\) to \(\nu(m)\).
They would block \(\nu\).
Thus every woman weakly
prefers her partner in
\(\nu\) to her partner
in \(\mu\), proving
receiver pessimism.

Two proposer-optimal matchings
assign each proposer the
same partner by strictness,
so are identical. Every
schedule satisfies the
proved properties and
therefore has that same
output. Interchanging
the sides proves the
receiver-optimal result.
The empty instance
makes zero proposals
and satisfies all
conclusions vacuously.
\end{proof}
\end{annexitem}

\begin{annexitem}{annex:p18.matching.assumption.boundaries}{thm:p18-matching-assumption-boundaries}{The exact three assumption failures}
\begin{proof}
With all partners tied, there are
no strict attacks. The entire
four-pair set is conflict-free
and has no outsiders, hence
is stable, but is not a
matching. With no acceptable
pair the argument universe
is empty, so its empty
set is stable; no acceptable
perfect matching exists.

In the infinite example,
the displayed \(S\) has
distinct rows and columns.
An unchosen pair with
man \(m_{i+1}\) is attacked
by his top-choice pair
\((m_{i+1},w_i)\).
Every pair \((m_0,w_i)\)
is attacked by that same
chosen pair through
\(w_i\)'s top preference.
These cases exhaust
outsiders, proving
stability despite the
unmatched \(m_0\).
The specified rankings
are strict total orders:
put the first choice
first and enumerate all
remaining indices once.
Thus the failed totality
is due to infinite
cardinality, not a tie
or an incomplete list.
\end{proof}
\end{annexitem}

\begin{annexitem}{annex:p18.roommate.argumentation}{thm:p18-roommate-argumentation}{Both directions for one-population partial matchings}
\begin{proof}
Two distinct pairs sharing a person
are strictly ordered by that person's
alternatives, so one attacks the
other. Disjoint pairs do not
attack, and no pair attacks
itself. Thus conflict-free sets
are precisely partial matchings.

For an unchosen pair \(\{a,b\}\),
only the selected pair containing
\(a\), or the selected pair
containing \(b\), can attack.
Each does so exactly when
the corresponding person prefers
the current partner to the
proposed one. If neither
attacks, strict completeness,
with every partner above
\(\bot\), says that both
people prefer the proposed
pair, so it blocks.
Conversely a blocking
pair has no selected
attacker. These are
the two stability
directions, including
unmatched participants.
Two unmatched people
would block by pairing
together, proving the
at-most-one bound.
The following example
proves the stated
cardinality boundary
for preferred extensions.
\end{proof}
\end{annexitem}

\begin{annexitem}{annex:p18.five.person.pairing}{thm:p18-five-person-pairing}{The complete five-person calculation and ordered-role counterinstance}
\begin{proof}
The pair \(e_0=\{m,w\}\) is
unattacked because both people
rank each other first. It
attacks every pair joining
either of them to a triangle
person. No admissible set
contains such a cross pair,
because it cannot defend
against the unattacked
\(e_0\).

The three triangle pairs
form the directed cycle
\[
 \{p_1,p_2\}\to\{p_1,p_3\}\to
                    \{p_2,p_3\}\to\{p_1,p_2\}.
\]
A conflict-free set contains
at most one of these
pairs. One selected pair
does not attack its
predecessor in the cycle;
\(e_0\) is disjoint from
all three, and cross
pairs have already been
excluded from admissible
sets. Therefore no
admissible set contains
a triangle pair.
Exactly \(\varnothing\)
and \(\{e_0\}\) remain.

The sole unattacked argument
is \(e_0\), and the same
attacker cases give
\(F(\{e_0\})=\{e_0\}\).
This is grounded and
uniquely preferred. It
does not attack the
triangle pairs, so
is not stable, and
no stable extension
exists. Two ordinary
disjoint pairs exist,
for example \(e_0\)
and \(\{p_1,p_2\}\);
three would require
six distinct people.
The maximum cardinality
is therefore two pairs,
not the preferred set's
one. None of these
arguments uses an
unspecified lower ranking.

For the ordered-role graph,
let \(f\) exchange \(m,w\)
and cycle the three
\(p_i\)'s as displayed.
The set \(S_f=\{(a,f(a)):a\in P\}\)
has distinct first and
second coordinates because
\(f\) is a permutation,
so is conflict-free.
Any outsider \((a,b)\)
has \(b\ne f(a)\).
The chosen \((a,f(a))\)
attacks it because
\(f(a)\) is \(a\)'s
first choice. This
also covers diagonals,
which were placed last.
Thus \(S_f\) is stable.
It has no diagonal
members, so removing
diagonals leaves it
stable. Each triangle
person appears with
two different partners
across the two roles;
it is not an
unordered matching.
\end{proof}
\end{annexitem}

\begin{annexitem}{annex:p18.matching.welfare.boundary}{thm:p18-matching-welfare-boundary}{The two stable matchings and stipulated utilities}
\begin{proof}
In \(\mu\), both proposers have
their first choices, so no
unchosen pair blocks. In
\(\nu\), both receivers have
their first choices, with
the same result. These
are the only two perfect
matchings of this instance.
Every proposer prefers
\(\mu\) and every receiver
prefers \(\nu\). Thus
\(\mu\) is proposer-optimal,
and neither matching
Pareto-dominates the other
for all participants.
The stipulated utility
sums are
\[
                       U(\mu)=1+1+0+0=2,\qquad
                       U(\nu)=0+0+10+10=20.
\]
Both numerical scales
represent exactly the
preference lists used
for stability. Hence
the stable proposer-optimal
\(\mu\) does not maximize
this sum. The interpersonal
weights were supplied
as additional data,
not inferred from
ordinal ranks.
\end{proof}
\end{annexitem}

\begin{annexitem}{annex:p18.superadditive.cover}{thm:p18-superadditive-cover}{The full partition-closure bridge}
\begin{proof}
For disjoint coalitions, combine their maximizing
partitions to obtain a partition of their union.
This gives
\(\widehat v(S\cup T)\geq\widehat v(S)+\widehat v(T)\).
The one-block partition gives
\(\widehat v(S)\geq v(S)\) for nonempty \(S\);
the empty case is equality. A singleton
has only its one-block partition, so
its value is unchanged. Together with
the stated equality at \(N\), these
facts make the two imputation sets
identical.

Any original domination remains feasible
under the larger values \(\widehat v\).
Conversely suppose \(x\) dominates \(y\)
via \(S\) for \(\widehat v\), and
choose a partition \(\pi\) attaining
\(\widehat v(S)\). If every block
\(K\in\pi\) had \(x(K)>v(K)\),
their sum would give
\(x(S)>\widehat v(S)\), contradicting
feasibility. Thus some nonempty
block has \(x(K)\leq v(K)\).
All its members still strictly
improve, since they belong to
\(S\). Therefore \(x\) already
dominates \(y\) via \(K\) for
the original function. The
argument graphs are literally
equal, proving preservation
of all their extension
families and NM solutions.
\end{proof}
\end{annexitem}

\begin{annexitem}{annex:p18.lucas.core.foundation}{thm:p18-lucas-core-foundation}{All characteristic values, core constraints and vertices}
\begin{proof}
The disjoint pairs \(Q_1,Q_2\) each have value
1, while their union has value zero, so
\(\ell\) is not superadditive. Every listed
coalition satisfies
\(\ell(S)\leq|S\cap O|\): check the five
classes in the definition, and use zero
for unlisted coalitions. Thus every
partition of \(N\) has total coalition
value at most \(|O|=5\). The one-block
partition attains \(\ell(N)=5\), proving
\(\widehat\ell(N)=5\). Singleton values
are zero and are preserved by the
cover. The preceding theorem gives
the same imputation set and domination
graph for the two functions.

Coalition inequalities for \(\ell\) and
for its cover are equivalent. Indeed,
if \(x(K)\geq\ell(K)\) for all \(K\),
then summing over any partition of
\(S\) gives \(x(S)\geq\widehat\ell(S)\).
The reverse follows from
\(\widehat\ell(S)\geq\ell(S)\).
Apply the proved core characterization
to the superadditive cover. Its
pair inequalities force \(x(Q_j)\geq1\);
their sum is 5, so all are equalities.
The \(O\) inequality forces \(x(O)\geq4\).
Thus a core point lies in the
displayed set.

Conversely take \(x\in B_\ell\) with
\(x(O)\geq4\). Each odd coordinate
is in \([0,1]\), and its paired
even coordinate is its complement
to 1. Removing one odd coordinate
from the odd sum leaves at least
3; removing two leaves at least
2. This verifies the three
odd four-member and six odd
three-member coalition inequalities.
For the mixed coalition \(\{1,4,7,9\}\),
put \(a=x_1,b=x_3,c=x_5,d=x_7,e=x_9\).
Its sum is
\[
 a+(1-b)+d+e
     =(a+b+c+d+e)+1-2b-c\geq5-2b-c\geq2.
\]
Cycling \(a,b,c\) verifies the
other two mixed coalitions.
Pairs and the grand coalition
already have their required
values, and all unlisted
inequalities follow from
nonnegativity. This proves
the exact core set.

For its convex representation put
\(\lambda_j=1-x_{2j-1}=x_{2j}\)
for \(1\leq j\leq5\) and
\(\lambda_0=1-\sum_j\lambda_j=x(O)-4\).
All weights are nonnegative
and sum to 1, and direct
coordinate comparison gives
\[
                       x=\lambda_0c^0+\sum_{j=1}^5\lambda_jc^{2j}.
\]
Conversely each listed vertex
is in the displayed core,
and that set is convex
by its linear constraints.
The two inclusions prove
the convex-hull formula.
No nonexistence conclusion
has been used in this
foundation proof.
\end{proof}
\end{annexitem}
